\section{Results}
\subsection{What the baseline exploits}
\label{sec:exploits}
A descriptor-only model comes close to StatNN. That model is ridge regression on hashed edit metadata (the family and parameter fields of each edit), with no weight statistics. On seed 0 of the trajectory split (345 rows, 278 clusters), the gap is 0.042 Spearman on MMLU-Pro and 0.035 on HellaSwag, computed from unrounded values (Table~\ref{tab:v2-baselines}, run \texttt{d3v2-m1-trajectory}). The descriptor also has the higher pooled out-of-fold $R^2$ on both targets: 0.560 against 0.555 on MMLU-Pro and 0.630 against 0.588 on HellaSwag. These seed-0 values come from \path{docs/council/2026-11-05/evidence-statnn-importances.md}, which \texttt{statnn\_evidence.py} generated, and they match the v2 seed-0 trajectory baselines on both targets.

On KL to base, which the v2 suite does not score, the pooled out-of-fold Spearman for seed-0 StatNN is 0.70, against 0.49 for the descriptor model, a lead of 0.20. Across three seeds, StatNN's paired MAE gain over the descriptor model is 4.1\% [$-$6.3\%, 14.0\%] on MMLU-Pro and 5.5\% [$-$6.3\%, 16.4\%] on HellaSwag, and both intervals include zero. These gains come from a separate metrology analysis (\path{docs/council/2026-11-05/round2_metrology_output.txt}, lines 100 and 102), not from a v2 run.

Table~\ref{tab:dynrange} shows how far each family moves the targets. Random-direction edits leave HellaSwag unchanged in all 60 rows, so their HellaSwag Spearman is undefined, and their MMLU-Pro SD is 0.007. Abliteration varies little, with SDs of 0.013 and 0.005. LoRA SFT has the largest mean shifts, $-$0.063 on MMLU-Pro and $+$0.034 on HellaSwag, and on both targets that mean exceeds the family's own SD. A pooled rank can therefore reflect family offsets as well as differences between edits.
\begin{table}[t]
\centering
\footnotesize
\begin{tabular}{lrll}\toprule
Family & Rows & MMLU-Pro $\Delta y$: mean (SD) [min, max] & HellaSwag $\Delta y$: mean (SD) [min, max] \\\midrule
Abliteration & 98 & 0.000 (0.013) [$-$0.040, 0.030] & 0.002 (0.005) [$-$0.010, 0.020] \\
Chain & 98 & $-$0.013 (0.038) [$-$0.140, 0.080] & 0.008 (0.016) [$-$0.030, 0.050] \\
LoRA SFT & 59 & $-$0.063 (0.049) [$-$0.190, 0.030] & 0.034 (0.013) [0.010, 0.070] \\
Random direction & 60 & 0.003 (0.007) [$-$0.010, 0.020] & 0.000 (0.000) [0.000, 0.000] \\
Task vector & 30 & $-$0.010 (0.034) [$-$0.110, 0.030] & 0.005 (0.015) [$-$0.020, 0.040] \\\bottomrule
\end{tabular}
\caption{Change in target accuracy per edit, $\Delta y$, from the matched predecessor, on the 345 seed-0 trajectory rows. Values are the mean, the sample SD (ddof 1), and the range. Rows are edits, not independent clusters. Source: \texttt{docs/council/2026-11-05/evidence-statnn-importances.md}, section ``Family dynamic range''.}
\label{tab:dynrange}
\end{table}

Table~\ref{tab:within} refits StatNN within each family, using five trajectory-grouped folds on seed 0. Within-family Spearman runs from 0.13 to 0.71 on MMLU-Pro and from $-$0.12 to 0.71 on HellaSwag. The 0.71 values come from the two smallest families: LoRA SFT (59 rows) on MMLU-Pro and task vector (30 rows) on HellaSwag. These refits have no intervals, and each uses a single family, so they show where rank varies within a family and are not predictors.
\begin{table}[t]
\centering
\footnotesize
\begin{tabular}{lrrrrr}\toprule
Family & Rows & MMLU-Pro & HellaSwag & KL to base & WikiText log-PPL \\\midrule
Abliteration & 98 & 0.27 & $-$0.12 & 0.50 & 0.19 \\
Chain & 98 & 0.46 & 0.51 & 0.45 & 0.35 \\
LoRA SFT & 59 & 0.71 & 0.01 & 0.89 & 0.69 \\
Random direction & 60 & 0.13 & undefined & 0.88 & $-$0.10 \\
Task vector & 30 & 0.22 & 0.71 & 0.82 & 0.68 \\\bottomrule
\end{tabular}
\caption{Within-family StatNN refits, Spearman, seed 0. Each family subset is refit alone with five trajectory-grouped folds. Source: \texttt{docs/council/2026-11-05/evidence-statnn-importances.md}, section ``Within-family refits''. The KL and WikiText columns are seed-0 analyses, not v2 suite outputs. Undefined means a constant observed or predicted value.}
\label{tab:within}
\end{table}

In StatNN's split-gain allocation on MMLU-Pro, the top five features hold 26\% of total gain. Three of those five are depth-bucket-0 statistics, from the earliest layers, of the down, o and v projections (source as above, section ``Gain importances''). These shares describe one fitted model. They are not causal importances.
